\section*{Appendix}

\subsection*{A. Software Environment and Reproduction}

Calculations use Python 3.12.13, with NumPy, pandas, Matplotlib, Seaborn, OpenPyXL, SciPy, and HiGHS for scientific computing, data processing, plotting, and optimization. The environment is managed by \texttt{uv}. Dependency constraints and locked versions are recorded in \texttt{pyproject.toml} and \texttt{uv.lock} in the supporting materials.

To reproduce results, first synchronize the environment, then execute the problem-level entry points in order from Problem 1 to Problem 4, working in the repository's \texttt{code/} directory. Preprocessing for Problems 2--4 depends on the shared processed layer generated by Problem 1. Run:

\begin{lstlisting}[language={},basicstyle=\ttfamily\small]
uv sync
uv run python question\question_01\main.py
uv run python question\question_02\main.py
uv run python question\question_03\main.py
uv run python question\question_04\main.py
\end{lstlisting}

Each entry point performs data processing, model solving, validation, and plotting in sequence. The supporting materials retain processed model inputs and generated result tables. To check published values, inspect result and independent-verification tables directly. To recompute from raw attachments, repeat preprocessing while preserving the original problem data.

\subsection*{B. Source Code and Supporting Materials}

Supporting materials include all source programs, environment configuration, processed model inputs, key result and validation tables for Problems 1--4, paper figures, and file checksum inventories. Result tables trace specific reported values; figures present results; independent-verification tables check job coverage, resource capacity, energy balance, storage states, and time boundaries.

Problem-specific evidence includes demand forecasting, baseline scheduling, and stress boundaries for Problem 1; multiobjective scheduling, baseline comparison, resource--energy recomputation, and sensitivity tests for Problem 2; storage optimization, regional metrics, baseline checks, and constraint tests for Problem 3; and joint computing--storage--electricity scheduling, scenarios, representative-window tests, and independent verification for Problem 4.

\subsection*{C. Limits on Result Interpretation}

Reported values come from the corresponding result tables, not estimates from figures. Passing hard-constraint checks establishes feasibility only for the given data, parameters, and boundaries. It does not establish global optimality or universal superiority across scenarios. Matheuristic results are reported with their actual solution method, comparison tests, and remaining limitations.

Hours 0--2399 form the main operating period; hours 2400--2405 complete already arrived jobs; hour 2406 is the terminal settlement boundary, not an ordinary execution hour. Costs including electricity sale revenue are net settlement values and do not imply negative actual electricity consumption.

\subsection*{D. AI Tool Use Details}

The team used DeepSeek during the competition for language editing, code debugging, figure preparation, and LaTeX typesetting assistance. AI was not used to determine core models, formulas, parameters, raw data, or final conclusions.

\subsubsection*{D.1 Tool and Dates of Use}

\begin{table}[H]
    \centering
    \small
    \renewcommand{\arraystretch}{1.05}
    \setlength{\tabcolsep}{3pt}
    \begin{tabularx}{0.96\textwidth}{>{\centering\arraybackslash}p{0.18\textwidth}>{\centering\arraybackslash}p{0.20\textwidth}>{\centering\arraybackslash}X>{\centering\arraybackslash}p{0.22\textwidth}}
        \toprule
        AI tool & Version/model & Developer & Dates of use \\
        \midrule
        DeepSeek & DeepSeek-v4-flash-0731 & DeepSeek & August 7--10, 2026 \\
        \bottomrule
    \end{tabularx}
\end{table}

The model version used was DeepSeek-v4-flash-0731.

\subsubsection*{D.2 Activities and Purposes}

DeepSeek assisted with language editing, debugging, figure preparation, final-draft compression, terminology consistency, citations, result-version checks, and LaTeX review. AI neither generated nor altered raw data and did not change the objectives, constraints, or parameters of the existing joint job search. Additional results were generated by actual local execution and independently recomputed.

\subsubsection*{D.3 Key Interaction Records}

\noindent\textbf{Record 1: DeepSeek language assistance, August 7--10, 2026}

The prompt requested review only of text written by the team, with local revision suggestions and no new models, formulas, data, or conclusions. AI identified long sentences, repetition, inconsistent terminology, and weak connections between text and figures. The team manually screened, checked, and rephrased suggestions.

\noindent\textbf{Record 2: Code debugging and figure assistance, August 9, 2026}

The prompt restricted review to syntax, indexing, data types, file reading, and plotting parameters in existing programs. It prohibited changes to raw data, objectives, constraints, model structure, parameters, and final conclusions. AI suggested variable checks, index corrections, data-type conversions, and figure-label improvements. The team adopted only suggestions reproducible locally without changing model meaning.

\noindent\textbf{Record 3: DeepSeek final-draft review, August 10, 2026}

The prompt required strict review of Problem 4 result accounting against a checklist, removal of mixed data versions across the sequential baseline, joint schedule, figures, and text, and consistent academic language, citations, and LaTeX formatting. AI assisted in adding entry points for consistent fixed-job energy recomputation and independent verification, and in revising plotting and paper sources. Every new metric was generated by actual execution on the team's computer; the team performed final checks and submission.

\subsubsection*{D.4 Adoption and Manual Revision}

The team checked all AI suggestions before adoption. Core methods, formulas, parameters, results, and conclusions were determined by the team, and final programs, data, figures, and conclusions were locally executed and verified. Complete DeepSeek records are contained in the corresponding use report in the supporting materials.

AI-assistance statements in the programs refer to the code checks and revisions involving DeepSeek. Assistance was limited to program review, result checking, figure preparation, and paper formatting; it did not rewrite the existing joint job-search model.

\subsection*{E. Source Inventory and Attachment Location}

All 24 Python source programs used for this problem were included in the source-code directory of the separately submitted support ZIP instead of being reprinted in the paper. That attachment also contains \texttt{pyproject.toml}, \texttt{uv.lock}, processed model inputs, key result and validation tables, paper figures, and a SHA-256 checksum inventory. Raw data supplied by the competition were not duplicated. The local readable copy is named \texttt{CCM2601033-support.zip}; the original competition package is retained in the frozen archive.

The source programs follow the relative directories for Problems 1--4 and comprise 24 \texttt{.py} files. Reproduction uses the commands in Appendix A in sequence. Paper and support attachments were submitted separately. The original attachment name followed the required problem-choice, entry-number, and attachment format, and its size was kept below 20 MB.
